
\section{Polynomical Chaos Expansions}
Polynomial chaos expansions (PCE) are a spectral method for uncertainty quantification. A spectral mathod is a method where a
stochastic process is represented as a series expansion using an orthogonal basis of functions. Spectral methods
decompose a random variable into components, similar to a fourier analysis \parencite{heuvelineHybridGeneralizedPolynomial2014}.

\subsection{Generalised Polynomial chaos}
PCE uses polynomials as the orthogonal basis. Originally developed to model Wiener processes using Hermite polynomials, it was extended by
\textcite{xiuWienerAskeyPolynomial2002} to use other polynomials and thus represent other probability distributions. The
generalised polynomial chaos expansion (gPC) works as follows:
\TODO{Fix bold greek symbols}
Let \(X\) be a random variable in probability space \(\left(\Omega, \mathcal{F}, \mathbb{P}\right)\), where \(\Omega\) is
the sample space with \(\omega \in \Omega\). \(X\) can then be formulated as follows
\parencite{heuvelineHybridGeneralizedPolynomial2014}:

\begin{equation}
\label{eq:gPC_explain_X}
X(\omega) = \sum_{i = 0}^{\infty}c_i\boldsymbol{\Psi}_i(\boldsymbol{\xi})
\end{equation}

\(\boldsymbol{\xi}\) is a random variable subject to some probability distribution.\({\boldsymbol{\Psi}}\) is a set of orthogonal polynomials that describes the
distribution of the random variable. \(c_i\) is a weight given to each term of the polynomial. These are determined by
pointwise sampling at a given point in time. Given a time dependent process \(X = X(t, \omega)\), \(c_i\) has to be
re-evaluated at each timestep. 

infinite sums are not possible in reality, so the sum is truncated at point \(M\). This gives an approximation of \(X\)
defined as:

\begin{equation}
X^M(t, \omega) \coloneq \sum_{i = 0}^{M}c_i(t)\boldsymbol{\Psi_i}(\boldsymbol{\xi})
\end{equation}
Where M is given by:
\begin{equation}
M = \frac{(P + L)!}{P!L!} - 1
\end{equation}

The truncation order \(M\) is dependent on two discretisation parameters, \(L\) and \(P\). \(P\) is the maximal polynomial order
used, and \(L\) is the dimension of the vector of random variables \(\boldsymbol{\xi} = (\xi_1, \dots, \xi_L)\). The value of M
grows rapidly as the dimensionality of the problem increases. The mean and variance of the result can be determined from
the resulting expansion. 

\begin{align}
    \overbar{X}(t) &= c_0(t) \\
    \mathit{Var} \, \left[ X(t) \right] &= \sum_{i = 1}^{M} c_i(t)^2
\end{align}

Overall this method works more efficiently than Monte-Carlo as the coefficients for the sum can be estimated using far
fewer model evaluations. Additionally, once the polynomial surrogate model is constructed, it can be evaluated quickly
for any variations of the input values. 

\subsection{Time dependent generalised polynomial chaos}
\textcite{gerritsmaTimedependentGeneralizedPolynomial2010} introduced the concept of time-dependent generalised
polynomial chaos (TD-gPC). One of the limitaions of gPC is a lack of convergence when the dependence on \(\boldsymbol{\xi}\) is
strongly non-linear. This is especially the case when considering long term integration
\parencite{heuvelineHybridGeneralizedPolynomial2014}. The difference between this method and the time independent method
is that with TD-gPC, the orthogonal polynomials can be adjusted if the distributions of \( \boldsymbol{\xi} \) changes.
This increases the accuracy of the PCE over long-term simulations. 

The TD-gPC approach does have some limitations in the application to dynamic systems. For longer simulations, the
accuracy of the method decreases \parencite{gerritsmaTimedependentGeneralizedPolynomial2010}, thus needing recalculation
and adding extra computing cost. \textcite{ebadollahiTimedependentUncertaintyQuantification2025} also mentions a
deterioration of the performance of TD-gPC for longer simulations. The computational cost rapidly increases as the
required polynomial order increases as time evolves. This makes it inefficient for strongly nonlinear dynamic systems.
The technique does not scale well for higher dimension problems
\parencites{heuvelineHybridGeneralizedPolynomial2014}{maiSurrogateModellingStochastic2016}. 


\subsection{Polynomial Dimensional Decomposition}
The poor scaling of PCE makes it less suitable for higher dimensional systems. An alternative approach, based on the
same mathematical basis, is polynomial dimensional decomposition (PDD). 
